\chapter{OpenGL I: The Independent Compatibility Ladder}
\label{ch:opengl-fixed-profiles}

CNA does not have one OpenGL renderer with a version switch.  It has several independent
families that happen to speak related APIs.  This chapter follows the compatibility ladder:
desktop OpenGL~1.x through \texttt{OPENGL1}, the GLSL~1.10 transition through
\texttt{OPENGL2}, and mobile fixed function through \texttt{OPENGLES1}.  Chapter~\ref{ch:opengl-programmable}
then covers the five-profile EasyGL family, \texttt{OPENGL4}, and \texttt{PORTABLEGL}.

That distinction matters operationally.  Selecting a newer identity does not mutate an older
implementation, and a fix in one family does not automatically reach another.  The renderer
contract is shared; context creation, resource ownership, draw submission, and evidence are
not.

\section{Three identities, three implementation families}

\begin{longtable}{p{0.16\textwidth}p{0.18\textwidth}p{0.22\textwidth}p{0.31\textwidth}}
\toprule
Identity & Pipeline & Shader delivery & Deliberate boundary \\
\midrule
\endfirsthead
\toprule
Identity & Pipeline & Shader delivery & Deliberate boundary \\
\midrule
\endhead
\texttt{OPENGL1} & desktop fixed function & none & OpenGL~1.x concepts; optional FBO and query entry points are detected at runtime \\
\texttt{OPENGL2} & programmable compatibility profile & GLSL~1.10 source compiled by the driver & compatibility-profile state plus shaders; MRT and queries depend on loaded functions \\
\texttt{OPENGLES1} & mobile fixed function & none & OpenGL ES~1.1; no MRT and no occlusion-query contract \\
\bottomrule
\end{longtable}

All three implement CNA directly.  None routes through EasyGL, and none is an alias for an
EasyGL profile.  The apparently similar names therefore describe API eras, not modes of one
code base.

\section{OPENGL1: a modern contract on a fixed pipeline}

\texttt{OPENGL1} is available only on Linux or Windows and links the system OpenGL library.
There is no shader stage to compile, cache, or bind.  CNA state is translated into texture
environment, matrix, lighting, fog, alpha, blend, depth, stencil, and client-array state that
the fixed pipeline understands.

This is not equivalent to a 2D fallback.  Both extended primitive entry points are native
overrides and submit fixed-function geometry.  A \cnaclass{GpuDrawParams} packet must be
reduced to the subset expressible by the pipeline rather than silently falling through the
base colored path.  The result can draw 3D geometry, but custom programmable effects are an
architectural impossibility, not a missing weekend task.

\subsection{Optional extensions are part of the public truth}

OpenGL~1.x predates framebuffer objects and standardized occlusion queries.  The renderer
therefore treats both as runtime discoveries:

\begin{itemize}
  \item \cnaclass{RenderTarget2D} is available only when the required FBO entry points were
        resolved.  Binding the default backbuffer remains valid without them.
  \item multiple render targets are not implemented;
  \item occlusion queries are available only when the ARB or later core query functions can
        be loaded.
\end{itemize}

This conditional shape is more honest than returning an unconditional capability bit based
only on the build identity.  An executable can be built successfully and still meet a
driver whose OpenGL~1.x extension surface is narrower.

\begin{sourcenote}
The platform gate and system-library dependency are build facts; the FBO and query clauses
are runtime facts.  A build matrix can prove the former, but only execution against the
selected driver can prove the latter.
\end{sourcenote}

\section{OPENGL2: the GLSL 1.10 bridge}

\texttt{OPENGL2} is the transition renderer.  Its context remains a compatibility-profile
world, while vertex and fragment work is expressed as GLSL~1.10.  That combination makes it
materially different from both neighbours: it is programmable unlike \texttt{OPENGL1}, but
cannot be served by EasyGL's ES-oriented shader rewriting and profile macros.

The renderer owns real extended primitive paths, render targets, and effect-aware submission.
Its optional ceiling again follows loaded functions rather than the name alone:

\begin{itemize}
  \item 2D render targets are supported;
  \item MRT is conditional on \texttt{glDrawBuffers};
  \item occlusion queries are conditional on query creation being available.
\end{itemize}

GLSL~1.10 also explains why shader source is an integration boundary.  A shader written only
for modern core GLSL cannot simply be handed to this renderer, and an ES shader is not made
desktop-compatible merely by changing its version line.  Attribute syntax, fragment outputs,
precision qualifiers, and built-ins belong to the selected language profile.

\subsection{Why keep this renderer beside EasyGL?}

Compatibility hardware and virtualized drivers can expose a usable OpenGL~2.x path without
providing the ES or modern-core environment expected elsewhere.  Keeping this implementation
separate preserves that deployment option and makes its compromises reviewable.  Folding it
into EasyGL would hide a different shader language and a different context contract behind a
single selector.

\section{OPENGLES1: fixed function on the mobile branch}

\texttt{OPENGLES1} targets \texttt{libGLESv1\_CM} and \texttt{GLES/gl.h}.  The dependency is
not vendored, and the module refuses configuration when the system library is absent.  Stock
Mesa builds commonly disable GLES1, so the repository supplies a dedicated ES1-enabled Mesa
environment script for this renderer's evidence path.

Like \texttt{OPENGL1}, it has no shaders.  Unlike \texttt{OPENGL1}, its API ceiling is the
smaller ES~1.1 fixed-function surface.  It implements 2D targets but neither MRT nor occlusion
queries.

\subsection{A hybrid draw route, not unconditional effect parity}

The extended draw methods are overridden, so classifying this renderer as a base fallback
would be wrong.  Their internal dispatch is nevertheless hybrid.  Combinations the fixed
pipeline can represent are handled natively; skinning, PBR, custom effects, instancing, and
other unavailable combinations are sent to the colored tail.  The call may therefore draw
geometry while losing the semantic content of \cnaclass{GpuDrawParams}.

This is narrower than the unconditional downgrade described for \texttt{DIRECTX10} in
Chapter~\ref{ch:renderer-contract}: supported fixed-function cases remain native here.  It is
still a crucial portability warning: successful submission does not prove that the requested
effect model survived dispatch.

\section{The same contract does not imply the same proof}

The three families require different verification environments.

\begin{longtable}{p{0.18\textwidth}p{0.34\textwidth}p{0.34\textwidth}}
\toprule
Identity & Strongest practical path & Remaining evidence boundary \\
\midrule
\endfirsthead
\toprule
Identity & Strongest practical path & Remaining evidence boundary \\
\midrule
\endhead
\texttt{OPENGL1} & native Linux or Windows system OpenGL, including extension-dependent cases & driver variation in FBO and query availability \\
\texttt{OPENGL2} & compatibility-profile desktop OpenGL with real shader compilation & loaded-function and GLSL-driver variation \\
\texttt{OPENGLES1} & custom Mesa build with GLES1 enabled & specialized environment is not the stock developer path \\
\bottomrule
\end{longtable}

Registration totals are deliberately omitted here.  They drift as cases are split or merged
and do not measure assertion quality.  The useful question is whether a test reaches a real
context, submits pixels through the family-specific path, and observes the result.

\section{Choosing among the compatibility identities}

Choose \texttt{OPENGL1} only when fixed-function desktop compatibility is itself a goal.
Choose \texttt{OPENGL2} when GLSL~1.10 and compatibility contexts match the deployment target.
Choose \texttt{OPENGLES1} when the ES~1.1 environment is intentional and its fixed-function
ceiling is acceptable.  For a broader programmable profile set, continue with
Chapter~\ref{ch:opengl-programmable}; a similar API surname does not erase the implementation
boundary.

Before shipping on any of the three, verify four things on the target driver:

\begin{enumerate}
  \item the selected identity reported at runtime is the one configured at build time;
  \item the render-target and query features actually used by the game are available;
  \item representative effect-aware draws retain their intended semantics rather than taking
        a reduced path;
  \item pixel evidence is collected from the real context, not only from a construction or
        no-throw smoke test.
\end{enumerate}

That checklist is the recurring lesson of the compatibility ladder: an old API can implement
a surprising amount of CNA, but its limits must be stated as limits of a particular path,
not blurred into the capabilities of the OpenGL family as a whole.
